\subsection{Patch-Quality Sensitivity}

We next evaluate whether performance improves when evaluation is restricted to
patches with stronger STHELAR quality control. We use the patch-level QC
Jaccard score as an external quality indicator and recompute metrics on
cumulative subsets with QC Jaccard above fixed thresholds. This analysis is not
used for model selection; instead, it quantifies how much of the residual error
is associated with lower-quality ST-derived supervision.

\begin{table}
\caption{KLT five-class QC sensitivity. Metrics are recomputed after retaining
test patches with QC Jaccard above the indicated threshold. Values are averaged
over seeds 42 and 43. \(\Delta\mPQ{}\) is the absolute gain over the full test
set for the same method.}
\label{tab:klt-qc}
\centering
{\fontsize{8}{9}\selectfont
\begin{tabular}{lcccccc}
\toprule
Method & QC subset & Retained & \mPQ{} & \(\Delta\mPQ{}\) & \bPQ{} & \Fdet{} \\
\midrule
\FullFT{} & All & 4653 (100.0\%) & 0.309 & -- & 0.515 & 0.829 \\
\FullFT{} & \(J_{\mathrm{QC}}\ge0.50\) & 3644 (78.3\%) & 0.363 & +0.054 & 0.593 & 0.850 \\
\FullFT{} & \(J_{\mathrm{QC}}\ge0.60\) & 2919 (62.7\%) & 0.386 & +0.077 & 0.626 & 0.857 \\
\FullFT{} & \(J_{\mathrm{QC}}\ge0.70\) & 1246 (26.8\%) & 0.431 & +0.121 & 0.674 & 0.873 \\
\midrule
\LoRA{}+\AF{}+heads & All & 4653 (100.0\%) & 0.294 & -- & 0.504 & 0.835 \\
\LoRA{}+\AF{}+heads & \(J_{\mathrm{QC}}\ge0.50\) & 3644 (78.3\%) & 0.343 & +0.049 & 0.584 & 0.858 \\
\LoRA{}+\AF{}+heads & \(J_{\mathrm{QC}}\ge0.60\) & 2919 (62.7\%) & 0.366 & +0.073 & 0.618 & 0.866 \\
\LoRA{}+\AF{}+heads & \(J_{\mathrm{QC}}\ge0.70\) & 1246 (26.8\%) & 0.410 & +0.117 & 0.667 & 0.883 \\
\bottomrule
\end{tabular}
}
\end{table}

Table~\ref{tab:klt-qc} shows that both \FullFT{} and the selected PEFT adapter
improve monotonically as the evaluation subset is restricted to higher-QC
patches. For \LoRA{}+\AF{}+heads, \mPQ{} increases from 0.294 on the full KLT
test set to 0.343 at \(J_{\mathrm{QC}}\ge0.50\), 0.366 at
\(J_{\mathrm{QC}}\ge0.60\) and 0.410 at \(J_{\mathrm{QC}}\ge0.70\).
\bPQ{} and detection F1 increase in parallel,
indicating that the effect is not limited to type assignment. The similar
absolute gains for \FullFT{} and PEFT suggest that the QC effect reflects the
evaluation target rather than a peculiarity of the adapter.

\begin{table}
\caption{Optional tissue-specific PEFT QC sensitivity. Values use the selected
\LoRA{}+\AF{}+heads tissue adapters reported in Table~\ref{tab:tissue}.
Very small high-QC subsets should be interpreted cautiously.}
\label{tab:tissue-qc}
\centering
{\fontsize{8}{9}\selectfont
\begin{tabular}{lccccc}
\toprule
Tissue & \mPQ{} & \multicolumn{2}{c}{\(J_{\mathrm{QC}}\ge0.60\)} &
\multicolumn{2}{c}{\(J_{\mathrm{QC}}\ge0.70\)} \\
\cmidrule(lr){3-4}\cmidrule(lr){5-6}
 & All & Retained & \mPQ{} & Retained & \mPQ{} \\
\midrule
Liver & 0.393 & 1992 (75.5\%) & 0.441 & 1012 (38.3\%) & 0.469 \\
Kidney & 0.230 & 482 (51.3\%) & 0.292 & 230 (24.5\%) & 0.297 \\
Ovary & 0.300 & 1703 (71.6\%) & 0.361 & 942 (39.6\%) & 0.396 \\
Breast & 0.335 & 5577 (59.1\%) & 0.397 & 1721 (18.2\%) & 0.456 \\
Colon & 0.205 & 834 (25.7\%) & 0.288 & 101 (3.1\%) & 0.360 \\
Lung & 0.302 & 1195 (48.4\%) & 0.361 & 339 (13.7\%) & 0.385 \\
Pancreatic & 0.239 & 465 (17.8\%) & 0.362 & 50 (1.9\%) & 0.426 \\
Skin & 0.224 & 752 (35.3\%) & 0.401 & 449 (21.1\%) & 0.422 \\
Tonsil & 0.235 & 2468 (57.2\%) & 0.311 & 678 (15.7\%) & 0.351 \\
\midrule
Mean & 0.273 & -- & 0.357 & -- & 0.396 \\
\bottomrule
\end{tabular}
}
\end{table}

The same pattern is visible across tissue-specific adapters
(Table~\ref{tab:tissue-qc}): mean PEFT \mPQ{} increases from 0.273 on all
patches to 0.357 at \(J_{\mathrm{QC}}\ge0.60\) and 0.396 at
\(J_{\mathrm{QC}}\ge0.70\). This supports the interpretation that a meaningful
part of the measured error comes from lower-confidence STHELAR patches. However,
strict thresholds can leave small subsets for some tissues, especially
pancreatic and colon, so these tissue-level QC values are best used as a
sensitivity analysis rather than as primary benchmark numbers.
